\documentclass[11pt]{paper}

\usepackage[T1]{fontenc}

\usepackage{graphicx}
\usepackage{amssymb}
\usepackage{amstext}
\usepackage{amsmath}
\usepackage{a4wide,color}
\usepackage[utf8]{inputenc}
\usepackage[frenchb]{babel}
\usepackage{xspace}
\usepackage{anysize}
\usepackage{tabularx}
\usepackage{multirow}
\usepackage{fancybox}
\usepackage{fancyhdr}
\usepackage{bbding}

\usepackage{color}
\usepackage{float}

\usepackage[toc,page]{appendix} 
\usepackage{lscape}

\usepackage{placeins}

\usepackage{listings}

\usepackage{todonotes}

\usepackage[framemethod=TikZ]{mdframed}

\lstset{% general command to set parameter(s)
	basicstyle=\footnotesize, % print whole listing small
	keywordstyle=\color{magenta}\bfseries, % underlined bold black keywords
	identifierstyle=, % nothing happens
	commentstyle=\color{black}, % white comments
	stringstyle=\ttfamily, % typewriter type for strings
	showstringspaces=false % no special string spaces
} 

\lstdefinelanguage{aald}
	{morekeywords={system, implementation},
	sensitive=false,
	morecomment=[l]{//},
	morecomment=[s]{/*}{*/},
	morestring=[b]",
}

\usepackage[colorlinks=true]{hyperref}

\renewcommand{\appendixtocname}{Annexes} 
\renewcommand{\appendixpagename}{Annexes}

\usepackage[normalem]{ulem}
\usepackage{color}

\definecolor{Fond}{gray}{0.7}

\renewcommand{\floatpagefraction}{.99}
\renewcommand{\textfraction}{.01}
\newcommand{\modif}[1]{\textcolor{red}{\uline{#1}}}

\newcommand{\cpp}{C{\tt ++}\xspace}

%%%%%%%%%%%%%%%%%%%%%%%%%%%%%%%%
\newcounter{obj}[section]\setcounter{obj}{0}
\renewcommand{\theobj}{\arabic{section}.\arabic{obj}}
\newenvironment{objectif}[2][]{%
\refstepcounter{obj}%
\ifstrempty{#1}%
{\mdfsetup{%
frametitle={%
\tikz[baseline=(current bounding box.east),outer sep=0pt]
\node[anchor=east,rectangle,fill=blue!20]
{\strut Objectif~\theobj};}}
}%
{\mdfsetup{%
frametitle={%
\tikz[baseline=(current bounding box.east),outer sep=0pt]
\node[anchor=east,rectangle,fill=blue!20]
{\strut (\theobj) Objectif~:~#1};}}%
}%
\mdfsetup{innertopmargin=10pt,linecolor=blue!20,%
linewidth=2pt,topline=true,%
frametitleaboveskip=\dimexpr-\ht\strutbox\relax
}
\begin{mdframed}[]\relax%
\label{#2}}{\end{mdframed}}
%%%%%%%%%%%%%%%%%%%%%%%%%%%%%%%%
%%%%%%%%%%%%%%%%%%%%%%%%%%%%%%%%
\newcounter{wtodo}[section]\setcounter{wtodo}{0}
\renewcommand{\thewtodo}{\arabic{section}.\arabic{wtodo}}
\newenvironment{wtodo}[2][]{%
\refstepcounter{wtodo}%
\ifstrempty{#1}%
{\mdfsetup{%
frametitle={%
\tikz[baseline=(current bounding box.east),outer sep=0pt]
\node[anchor=east,rectangle,fill=red!20]
{\strut Travail à faire~\thewtodo};}}
}%
{\mdfsetup{%
frametitle={%
\tikz[baseline=(current bounding box.east),outer sep=0pt]
\node[anchor=east,rectangle,fill=red!20]
{\strut (\thewtodo) Travail à faire~:~#1};}}%
}%
\mdfsetup{innertopmargin=10pt,linecolor=red!20,%
linewidth=2pt,topline=true,%
frametitleaboveskip=\dimexpr-\ht\strutbox\relax
}
\begin{mdframed}[]\relax%
\label{#2}}{\end{mdframed}}
%%%%%%%%%%%%%%%%%%%%%%%%%%%%%%%%
%%%%%%%%%%%%%%%%%%%%%%%%%%%%%%%%
\newcounter{howto}[section]\setcounter{howto}{0}
\renewcommand{\thehowto}{\arabic{section}.\arabic{howto}}
\newenvironment{howto}[2][]{%
\refstepcounter{howto}%
\ifstrempty{#1}%
{\mdfsetup{%
frametitle={%
\tikz[baseline=(current bounding box.east),outer sep=0pt]
\node[anchor=east,rectangle,fill=green!20]
{\strut Comment faire~\thehowto};}}
}%
{\mdfsetup{%
frametitle={%
\tikz[baseline=(current bounding box.east),outer sep=0pt]
\node[anchor=east,rectangle,fill=green!20]
{\strut (\thehowto) Comment faire~:~#1};}}%
}%
\mdfsetup{innertopmargin=10pt,linecolor=green!20,%
linewidth=2pt,topline=true,%
frametitleaboveskip=\dimexpr-\ht\strutbox\relax
}
\begin{mdframed}[]\relax%
\label{#2}}{\end{mdframed}}
%%%%%%%%%%%%%%%%%%%%%%%%%%%%%%%%
%%%%%%%%%%%%%%%%%%%%%%%%%%%%%%%%
\newcounter{entracte}[section]\setcounter{entracte}{0}
\renewcommand{\theentracte}{\arabic{section}.\arabic{entracte}}
\newenvironment{entracte}[2][]{%
\refstepcounter{entracte}%
\ifstrempty{#1}%
{\mdfsetup{%
frametitle={%
\tikz[baseline=(current bounding box.east),outer sep=0pt]
\node[anchor=east,rectangle,fill=black!20]
{\strut Entracte~\theentracte};}}
}%
{\mdfsetup{%
frametitle={%
\tikz[baseline=(current bounding box.east),outer sep=0pt]
\node[anchor=east,rectangle,fill=black!20]
{\strut (\theentracte) Entracte~:~#1};}}%
}%
\mdfsetup{innertopmargin=10pt,linecolor=black!20,%
linewidth=2pt,topline=true,%
frametitleaboveskip=\dimexpr-\ht\strutbox\relax
}
\begin{mdframed}[]\relax%
\label{#2}}{\end{mdframed}}
%%%%%%%%%%%%%%%%%%%%%%%%%%%%%%%%
%%%%%%%%%%%%%%%%%%%%%%%%%%%%%%%%
\newcounter{bp}[section]\setcounter{bp}{0}
\renewcommand{\thebp}{\arabic{section}.\arabic{bp}}
\newenvironment{bonnepratique}[2][]{%
\refstepcounter{bp}%
\ifstrempty{#1}%
{\mdfsetup{%
frametitle={%
\tikz[baseline=(current bounding box.east),outer sep=0pt]
\node[anchor=east,rectangle,fill=orange!20]
{\strut Bonne pratique~\thebp};}}
}%
{\mdfsetup{%
frametitle={%
\tikz[baseline=(current bounding box.east),outer sep=0pt]
\node[anchor=east,rectangle,fill=orange!20]
{\strut (\thebp) Bonne pratique~:~#1};}}%
}%
\mdfsetup{innertopmargin=10pt,linecolor=orange!20,%
linewidth=2pt,topline=true,%
frametitleaboveskip=\dimexpr-\ht\strutbox\relax
}
\begin{mdframed}[]\relax%
\label{#2}}{\end{mdframed}}
%%%%%%%%%%%%%%%%%%%%%%%%%%%%%%%%
%%%%%%%%%%%%%%%%%%%%%%%%%%%%%%%%
\newcounter{ad}[section]\setcounter{ad}{0}
%\renewcommand{\theentracte}{\arabic{section}.\arabic{entracte}}
\newenvironment{ad}[2][]{%
\refstepcounter{ad}%
\ifstrempty{#1}%
{\mdfsetup{%
frametitle={%
\tikz[baseline=(current bounding box.east),outer sep=0pt]
\node[anchor=east,rectangle,fill=black!20]
{\strut Message à caractère informatif};}}
}%
{\mdfsetup{%
frametitle={%
\tikz[baseline=(current bounding box.east),outer sep=0pt]
\node[anchor=east,rectangle,fill=black!20]
{\strut (\theentracte) Entracte};}}%
}%
\mdfsetup{innertopmargin=10pt,linecolor=red!60,%
linewidth=2pt,topline=true,%
frametitleaboveskip=\dimexpr-\ht\strutbox\relax
}
\begin{mdframed}[]\relax%
\label{#2}}{\end{mdframed}}
%%%%%%%%%%%%%%%%%%%%%%%%%%%%%%%%


\pagestyle{fancy}
\fancyhf{}
\fancyhead[RE,RO]{\thepage}
\fancyhead[LE]{}
\fancyhead[LO]{GPGPU}
\fancyfoot[LO]{Centrale Nantes -- \today}
\fancypagestyle{plain}{%
  \fancyhf{} % get rid of headers
  \renewcommand{\headrulewidth}{0pt} % and the line
}

\newenvironment{maliste}%
{ \begin{list}%
	{\ArrowBoldRightStrobe}%
	{\setlength{\labelwidth}{30pt}%
	 \setlength{\leftmargin}{35pt}%
	 \setlength{\itemsep}{\parsep}}}%
{ \end{list} }



\title{{\Large Programmation sur Processeur Graphique -- GPGPU }\\
\vspace{-0.4em}
{\large TD 3 : kernel multidimension}\\
{\small Centrale Nantes}\\
\small P.-E. Hladik, \small{\texttt{pehladik@ec-nantes.fr}}\\
---\\
{\scriptsize  Version bêta (\today)}\\
}


\begin{document}

\maketitle

%\begin{table}[htdp]
%\caption{Suivi des modifications}
%\begin{center}
%\begin{tabular}{|c|c|c|p{8cm}|}
%\hline
%Date & Version & Auteur & Modifications\\
%\hline
%JJ/MM/12 & 0.3.X & P.-E. Hladik&  \\
%\hline
%\end{tabular}
%\end{center}
%\label{default}
%\end{table}%

%%%%%%%%%%%%%%%%%%%%%%%%%
  \section{Multiplication matricielle basique}
%%%%%%%%%%%%%%%%%%%%%%%%%

\begin{objectif}[]{}
 \begin{itemize}
	\item implémenter une routine de base de multiplication de matrices denses
\end{itemize}
\end{objectif}

\begin{wtodo}[Multiplication matricielle]{}

Le code disponible dans le fichier {\tt matrix.cu} fournit les fonctions suivantes :
\begin{enumerate}
	\item {\tt initMatrix} pour allouer des valeurs à une matrice,
	\item {\tt compareMatrix} pour comparer deux matrices,
	\item {\tt testMatrixMulCPU} pour tester le calcul sur CPU avec des matrices connues,
	\item {\tt testMatrixMulGPU} pour tester le calcul sur GPU avec des matrices connues.\\
\end{enumerate}

Les fonction {\tt computeMatrixMulCPU} et {\tt computeMatrixMulGPU} sont simplement prototypées mais ne sont pas implémetées.\\

Au lancement du programme, on peut passer comme paramètre les dimension de la matrice A et B, par exemple si l'exécutable est nommé {\tt matrix}:

{\tt \$ ./matrix 10 30 30 40}

avec 10 le nombre de lignes de A, 30 le nombre de colonnes de A et de lignes de B et 40 le nombre de colonnes de B. 
%Si le nombre de colonnes de A et de lignes de B sont différentes une exception est lancée (le {\tt assert} ligne 19 du code). 
Vous pouvez ne pas renseigner ces valeurs, elles sont fixées à 100, 200, 200 et 400 par défaut.\\

Ajoutez les éléments nécessaires dans le code pour :
\begin{itemize}
	\item Implémenter la fonction {\tt computeMatrixMulCPU} pour réaliser la multiplication matricielle sur le CPU,
	\item Implémenter le kernel pour le calcul sur GPU (l'allocation de la mémoire est faite),
	\item Mesurer le temps d'exécution de {\tt computeMatrixMulCPU} (avec gprof, voir TD 2.b) et de  {\tt computeMatrixMulGPU} (avec nvprof, voir TD 2.b),
	\item Que donne la fonction {\tt compareMatrix} pour comparer les résultats produits sur le CPU et sur le GPU ? Expliquez ce qui vous arrive.
\end{itemize}
\end{wtodo}


%%%%%%%%%%%%%%%%%%%%%%%%%
  \section{Ensemble de Julia en CUDA}
%%%%%%%%%%%%%%%%%%%%%%%%%

\begin{objectif}[]{}
 \begin{itemize}
	\item réaliser un kernel multidimension
\end{itemize}
\end{objectif}

\begin{wtodo}[Ensemble de Julia]{}
Les ensembles de Julia est un exemple d'ensembles fractals. Considérant deux nombres complexes, $c$ et $z_0$, l'ensemble de Julia correspondant est la frontière de l'ensemble des valeurs initiales $z_0$ pour lesquelles la suite 
$${ z_{n+1}=z_{n}^{2}+c}$$
est bornée.

Nous allons représenter cet ensemble par une image où chaque pixel sera considéré comme un nombre complexe $z_0$ dans un plan où 0 est le centre de l'image. Si la suite $(z_n)$ est bornée pour le pixel $z_0$ le pixel prend la couleur rouge, sinon il prend la couleur noir.

Un code séquentiel permettant de représenter l'ensemble de Julia est déjà produit (pour $c = -0.8 + 0.156i$. Le but est de passer ce code sur un GPU.

\begin{enumerate}
\item Récupérer l'archive TD3.
\item Compilez le code avec {\tt nvcc -o julia julia\_bmp.cu}
\item Exécutez le code et observez la belle image produite.
\item Modifier le code pour l'exécuter sur le GPU et observez les performances !
\end{enumerate}
\end{wtodo}

\begin{howto}[{Données}]{}
Avant toute chose il faut penser à :
\begin{itemize}
\item allouer les espaces mémoire de travail sur le {\it device}, 
\item faire les copies de l'{\it host} vers le {\it device} (si besoin),
\item récupérer à la fin les données du {\it device} vers l'{\it host}.
\end{itemize}
\end{howto} 

\begin{howto}[{Nombre de grids et taille}]{}
Dans un premier temps, pour ne pas vous embêter avec les test sur les bornes de l'image, prenez un nombre de grids et de thread multiple de 1000*1000 (dimension de l'image).
\end{howto} 

\begin{howto}[{\tt \_\_device\_\_}]{}
N'oubliez pas que pour appeler une fonction sur le GPU il faut qu'elle soit déclarée comme {\tt \_\_device\_\_}, même les méthodes d'un objet...
\end{howto} 

%%%%%%%%%%%%%%%%%%%%%%%%%
  \section{Retour sur l'ensemble de Julia (pour ceux qui sont en avance)}
%%%%%%%%%%%%%%%%%%%%%%%%%

\begin{objectif}[]{}
 \begin{itemize}
	\item avoir un rendu de l'ensemble de Julia plus beau
\end{itemize}
\end{objectif}

\begin{wtodo}[Coloriser l'ensemble de Julia]{}

Au lieu d'afficher un point rouge, essayer de trouver une coloration qui prend en compte le nombre d'itération nécessaire pour tester si le point est dans l'ensemble de Julia ou non.

\end{wtodo}



\end{document}